% =============================================================
% 6x9 Layout
% =============================================================
\usepackage[
  paperwidth=6in,
  paperheight=9in,
  left=0.75in,
  right=0.75in,
  top=0.6in,
  bottom=1.0in,
  heightrounded
]{geometry}

\raggedbottom

% =============================================================
% Memoir-ish paragraph + chapter behavior
% =============================================================
\setlength{\parindent}{0pt}
\setlength{\parskip}{1em}

\renewcommand{\chapnamefont}{}
\renewcommand{\printchaptername}{}
\renewcommand{\chapnumfont}{}
\renewcommand{\printchapternum}{}
\setlength{\beforechapskip}{0em}
\setlength{\afterchapskip}{2em}
\setsecnumdepth{none}

% =============================================================
% Contents formatting
% =============================================================
\usepackage{titlesec}

\makeatletter
\renewcommand\tableofcontents{%
  \begingroup
    \setlength{\parskip}{1pt}%
    \setlength{\parindent}{0pt}%
    \vspace{5em}%
    {\LARGE \bfseries Contents\par}%
    \vspace{2em}%
    \@starttoc{toc}%
  \endgroup
}
\makeatother

% =============================================================
% Heading spacing (Pandoc: ## -> \section, ### -> \subsection, etc.)
% =============================================================
\titlespacing*{\section}{0pt}{0.8em}{0.1em}      % ##
\titlespacing*{\subsection}{0pt}{0.7em}{0em}     % ###

\titleformat{\subsubsection}
  {\normalfont\normalsize\bfseries}{\thesubsubsection}{1em}{}          % ####
\titlespacing*{\subsubsection}{0pt}{0.5em}{0em}

\titleformat{\paragraph}
  {\normalfont\normalsize\bfseries}{\theparagraph}{1em}{}              % #####
\titlespacing*{\paragraph}{0pt}{0.3em}{0em}

\titleformat{\subparagraph}
  {\normalfont\normalsize\bfseries\itshape}{\thesubparagraph}{1em}{}   % ######
\titlespacing*{\subparagraph}{0pt}{0.2em}{0em}

% =============================================================
% Header/Footer
% =============================================================
\usepackage{fancyhdr}
\pagestyle{fancy}
\fancyhf{}
\fancyfoot[LE,RO]{\thepage}
\renewcommand{\headrulewidth}{0pt}
\renewcommand{\footrulewidth}{0pt}

% =============================================================
% Horizontal rule helper
% =============================================================
\newcommand{\myhrule}{%
  \par\vspace*{1em}%
  {\centering\rule{0.6\linewidth}{0.4pt}\par}%
  \vspace*{1em}%
}

% =============================================================
% Fonts (Latin + CJK)
% =============================================================
\usepackage{fontspec}
\usepackage{xeCJK}

\setmainfont{Noto Serif}[Ligatures=TeX]
\setCJKmainfont{Noto Sans CJK SC}

% Keep “smart quotes” etc. from being treated as CJK
\AtBeginDocument{%
  \XeTeXcharclass`‘=0
  \XeTeXcharclass`’=0
  \XeTeXcharclass`“=0
  \XeTeXcharclass`”=0
  \XeTeXcharclass`…=0
}

% =============================================================
% Unicode symbol coverage (Option A)
% - Goal: arrows / misc symbols from LLM output render instead of tofu.
% - We only auto-switch for the Symbols class so xeCJK stays in charge of CJK.
% =============================================================
\usepackage[Symbols]{ucharclasses}

% Choose a symbols font that likely exists in your container; fall back safely.
\newfontfamily\SymbolsFont{Noto Sans Symbols 2}[Scale=MatchLowercase]
\IfFontExistsTF{Noto Sans Symbols 2}{}{
  \IfFontExistsTF{Noto Sans Symbols2}{
    \renewcommand\SymbolsFont{\newfontfamily\SymbolsFont{Noto Sans Symbols2}[Scale=MatchLowercase]}
  }{
    \IfFontExistsTF{Noto Sans Symbols}{
      \renewcommand\SymbolsFont{\newfontfamily\SymbolsFont{Noto Sans Symbols}[Scale=MatchLowercase]}
    }{
      \renewcommand\SymbolsFont{\newfontfamily\SymbolsFont{DejaVu Sans}[Scale=MatchLowercase]}
    }
  }
}

\setDefaultTransitions{\begingroup}{\endgroup}
\setTransitionsForSymbols{\begingroup\SymbolsFont}{\endgroup}

% =============================================================
% Math-ish Unicode normalization (minimal, safe)
% =============================================================
\usepackage{unicode-math}
\IfFontExistsTF{STIX Two Math}{
  \setmathfont{STIX Two Math}
}{
  \setmathfont{Latin Modern Math}
}

\usepackage{newunicodechar}
\newunicodechar{÷}{\ensuremath{\div}}
\newunicodechar{×}{\ensuremath{\times}}
\newunicodechar{−}{-}
\newunicodechar{≤}{\ensuremath{\le}}
\newunicodechar{≥}{\ensuremath{\ge}}
\newunicodechar{≠}{\ensuremath{\ne}}
\newunicodechar{≈}{\ensuremath{\approx}}
\newunicodechar{…}{\ldots}

% =============================================================
% Lists (NO \item patch — avoids recursion / TeX capacity blowups)
% =============================================================
\usepackage{enumitem}

% Labels: use glyphs that will exist (math glyphs are very reliable)
\setlist[itemize,1]{label=\textbullet}
\setlist[itemize,2]{label=$\circ$}
\setlist[itemize,3]{label=$\triangleright$}

% Spacing/indent: tune here instead of redefining \item
\setlist[itemize,enumerate]{leftmargin=1.5em}
\setlist[itemize]{topsep=0em, parsep=0pt, itemsep=0.191em, partopsep=0pt}
\setlist[enumerate]{topsep=0em, parsep=0pt, itemsep=0.191em, partopsep=0pt}

% If you really want an “after list” gap similar to your old after=\vspace{...}
% prefer the list-level knob:
\setlist[itemize,enumerate]{after=\vspace{0.25em}}

% =============================================================
% Block Quote Styling
% =============================================================
\usepackage{xcolor}
\usepackage{mdframed}

\renewenvironment{quote}
  {\begin{samepage}
   \begin{mdframed}[
    backgroundcolor=gray!10,
    linecolor=gray!60,
    leftline=true, rightline=false, topline=false, bottomline=false,
    innerleftmargin=1em, innerrightmargin=1.618em,
    innertopmargin=1em, innerbottommargin=1em,
    skipabove=1em, skipbelow=0em
   ]}
  {\end{mdframed}\vspace{-0.5em}\end{samepage}}

% =============================================================
% Captions: hide labels
% =============================================================
\usepackage{caption}
\captionsetup[figure]{labelformat=empty}
\captionsetup[table]{labelformat=empty}

% =============================================================
% Plotting
% =============================================================
\usepackage{pgfplots}
\pgfplotsset{width=10cm,compat=1.9}
\usepgfplotslibrary{statistics}
